\documentclass[11pt,a4paper]{article}
\usepackage[T1]{fontenc}
\usepackage[utf8]{inputenc}
\usepackage{amsmath,amssymb,amsthm,mathtools}
\usepackage{booktabs,longtable,array}
\usepackage[margin=23mm]{geometry}
\usepackage{xcolor}
\usepackage[hidelinks]{hyperref}

\newtheorem{theorem}{Théorème}[section]
\newtheorem{proposition}[theorem]{Proposition}
\newtheorem{lemma}[theorem]{Lemme}
\newtheorem{corollary}[theorem]{Corollaire}
\theoremstyle{definition}
\newtheorem{definition}[theorem]{Définition}
\newtheorem{remark}[theorem]{Remarque}

\newcommand{\Rev}{\operatorname{rev}}
\newcommand{\In}{\operatorname{In}}
\newcommand{\diag}{\operatorname{diag}}
\newcommand{\T}{\mathsf{T}}
\newcommand{\NT}{\mathsf{NT}}

\title{Formalisation finie du morphisme analytique de réversion,\\
de sa tour naturelle et du transfert ternaire d'inertie\\
\large Extension CORE7Q\_A1BH, degrés 58 à 64}
\author{Audit formel fini CORE7Q}
\date{9 août 2026}

\begin{document}
\maketitle

\begin{abstract}
Ce document formalise trois objets qui doivent rester strictement typés. Le premier est
l'entrelaceur linéaire associé à une étoile finie munie d'une involution. Le deuxième est
la naturalité abstraite de la famille de ces entrelaceurs entre degrés successifs. Le
troisième est un transfert dans le monoïde des inerties, observé aux degrés
$58,60,62,64$. Les identités décimales $329$, $159$ et $59\leftrightarrow95$ sont
également établies comme identités de mots de chiffres, puis séparées explicitement des
preuves analytiques. Toutes les conclusions sont finies; aucune revendication RH, aucune
positivité globale et aucune fermeture du complément infini n'en résultent.
\end{abstract}

\section{Étoile réversible et entrelaceur entier}

\begin{definition}[Étoile de niveau $n$]
Soit $S_n$ un ensemble de cardinal $2n-1$ muni d'une involution $J_n$ possédant
exactement un point fixe $f$ et $n-1$ orbites doubles
$\{a_r,b_r\}$, $1\le r\le n-1$. On note $\mathbb{R}^{S_n}$ l'espace des
coordonnées indexées par $S_n$.
\end{definition}

Dans la base canonique $(e_s)_{s\in S_n}$, on définit les colonnes
\[
 u_0=e_f,\qquad u_r=e_{a_r}+e_{b_r},\qquad
 v_r=e_{a_r}-e_{b_r}.
\]
L'application $\Phi_n$ est la matrice entière dont les colonnes sont
$u_0,u_1,\ldots,u_{n-1},v_1,\ldots,v_{n-1}$.

\begin{theorem}[Entrelacement fini]
Pour tout $n\ge2$,
\[
 J_n\Phi_n=\Phi_n\diag(I_n,-I_{n-1}),
 \qquad |\det\Phi_n|=2^{n-1}.
\]
En particulier,
\[
 \dim E_+(J_n)=n,\qquad \dim E_-(J_n)=n-1.
\]
\end{theorem}

\begin{proof}
Le point fixe fournit une colonne propre de valeur propre $+1$. Sur chaque orbite
double, la somme est propre de valeur $+1$ et la différence de valeur $-1$. Après
permutation des lignes et colonnes, $\Phi_n$ est somme directe d'un bloc $[1]$ et de
$n-1$ blocs $\left[\begin{smallmatrix}1&1\\1&-1\end{smallmatrix}\right]$, de
déterminant $-2$. L'identité et le déterminant en découlent.
\end{proof}

Aux niveaux testés, la taille $2n-1$ et la scission sont
\[
59=30+29,\quad61=31+30,\quad63=32+31,\quad65=33+32.
\]

\section{Transport des formes et inertie}

Soient $G_n$ une forme de Gram positive et $Q_n$ une forme symétrique sur
$\mathbb{R}^{S_n}$. Pour un seuil commun $\mu$, posons
$L_n=Q_n-\mu G_n$. Le transport par congruence est
\[
 \widehat G_n=\Phi_n^{\!T}G_n\Phi_n,
 \quad \widehat Q_n=\Phi_n^{\!T}Q_n\Phi_n,
 \quad \widehat L_n=\widehat Q_n-\mu\widehat G_n.
\]

\begin{proposition}[Covariance et inertie]
Si $\Phi_n$ est inversible, alors
\[
 \In(G_n)=\In(\widehat G_n),\quad
 \In(Q_n)=\In(\widehat Q_n),\quad
 \In(L_n)=\In(\widehat L_n).
\]
Cette conclusion est une covariance par congruence. Elle n'implique ni
$J_n^TQ_nJ_n=Q_n$, ni une permutation des entrées brutes de $Q_n$.
\end{proposition}

\begin{proof}
Il s'agit de la loi d'inertie de Sylvester. La dernière phrase distingue une congruence
par un entrelaceur inversible d'un automorphisme brut, obligation algébrique plus forte
et non certifiée ici.
\end{proof}

\section{Tour naturelle abstraite}

Soit $\iota_n$ une inclusion abstraite d'une étoile de taille $2n-1$ dans une étoile
de taille $2n+1$. On définit
\[
 \kappa_n=\Phi_{n+1}\iota_n\Phi_n^{-1}.
\]

\begin{theorem}[Carré commutatif transporté]
La famille $(\Phi_n)$ satisfait exactement
\[
 \kappa_n\Phi_n=\Phi_{n+1}\iota_n.
\]
Si $\iota_n$ commute avec les involutions abstraites, alors
\[
 J_{n+1}\kappa_n=\kappa_nJ_n.
\]
\end{theorem}

\begin{proof}
La première identité résulte de la définition de $\kappa_n$. Pour la seconde, on
insère les identités d'entrelacement de $\Phi_n$ et $\Phi_{n+1}$ et la commutation de
$\iota_n$ avec les décompositions $E_+\oplus E_-$. Toutes les opérations sont finies.
\end{proof}

Cette naturalité abstraite est certifiée. Elle ne choisit pas un ancrage canonique
parmi les supports résiduels bruts. Les candidats observés sont
\[
\begin{array}{c|c|c}
\text{degré}&\text{nombre}&\text{orbites candidates}\\
\hline
58&1&(3,26)\\
60&2&(3,27),(6,24)\\
62&3&(5,26),(6,25),(8,23)\\
64&1&(3,29).
\end{array}
\]
L'injection brute conservant l'indice gauche envoie $(a,b)$ sur $(a,b+1)$.
Elle possède un bord $58\to60$ et un bord $60\to62$, mais aucun bord
$62\to64$. Il n'existe donc pas de fil concret sur les quatre niveaux pour cette
injection. L'unicité locale au degré $64$ ne répare pas cette rupture globale.

\section{Audit des identités décimales}

\begin{definition}[Réversion de largeur fixe]
Pour un entier $x$ écrit sur exactement $k$ chiffres, $\Rev_k(x)$ est l'entier obtenu
en renversant le mot de chiffres, avec conservation préalable des zéros initiaux.
Sur les mots, $\Rev(uv)=\Rev(v)\Rev(u)$; c'est une anti-involution du monoïde libre.
\end{definition}

\subsection{Identité 329}

Les égalités suivantes sont exactes :
\[
 \Rev_2(32)=23,\qquad 32-23=9,
\]
\[
 \boxed{329=10\cdot32+(32-\Rev_2(32))=11\cdot32-\Rev_2(32)},
\]
et
\[
 \Rev_3(329)=923=100\cdot9+23.
\]
Au degré $58$, $329$ est aussi la cardinalité de la clôture réversible minimale du
support résiduel étudié. Cette coïncidence est donc une identité décimale attachée à
une cardinalité certifiée.

Elle ne se généralise pas aux successeurs :
\[
\begin{array}{c|r|r|r|r|c}
\text{degré}&C&\Rev_3(C)&p&\Rev_2(p)&C\bmod10=p-\Rev_2(p)\\
\hline
58&329&923&32&23&\T\\
60&451&154&45&54&\NT\\
62&525&525&52&25&\NT\\
64&506&605&50&05&\NT.
\end{array}
\]
Ainsi, la règle $329$ n'est pas une loi de la tour des clôtures.

\subsection{Identité 159 et réversion 59--95}

Le mot $159$ code le couple ``un point fixe, étoile de taille $59$'' :
\[
159=\operatorname{concat}(1,59),\qquad \Rev_2(59)=95,
\]
\[
\boxed{\Rev_3(159)=951=\operatorname{concat}(95,1)}.
\]
Les successeurs du même encodage sont
\[
\begin{array}{c|c|c|c}
\text{degré}&1\Vert(2n-1)&\Rev_2(2n-1)&\Rev_3(1\Vert(2n-1))\\
\hline
58&159&95&951\\
60&161&16&161\\
62&163&36&361\\
64&165&56&561.
\end{array}
\]
Il s'agit d'une anti-involution exacte de mots de chiffres. Une cardinalité, même
écrite de cette manière, ne définit ni une application linéaire, ni une congruence de
formes, ni un entrelaceur analytique. Un tel objet exige des coefficients et une
identité commutative vérifiés séparément.

\section{Transfert ternaire d'inertie}

\begin{definition}[Monoïde des inerties]
On note $\mathcal I=\mathbb N^3$ l'ensemble des triplets $(p,q,z)$, avec addition
composante par composante. Le rang est $r(p,q,z)=p+q+z$.
\end{definition}

Pour $k\le p$, définissons la translation partielle
\[
 \tau_k(p,q,z)=(p-k,q+k,z).
\]

\begin{proposition}[Générateur ternaire exact]
Sur chaque parité, le passage de $Q$ à $Q-\mu G$ applique
\[
 g=(-1,+1,0).
\]
Sur la somme directe des deux parités, il applique
\[
 2g=(-2,+2,0),\qquad -2g=(+2,-2,0)
\]
pour le passage inverse. Ces vecteurs conservent le rang.
\end{proposition}

\begin{proof}
Chaque canal possède exactement une direction négative sous le même seuil $\mu$.
Une direction passe donc du compte positif au compte négatif dans chaque canal.
L'addition des inerties des deux canaux donne le facteur deux. La somme des composantes
du générateur est nulle, ce qui préserve le rang.
\end{proof}

Les certifications finies donnent :
\[
\begin{array}{c|c|c|c}
\text{degré}&\In(Q)&\In(Q-\mu G)&\Delta\\
\hline
58&(59,0,0)&(57,2,0)&(-2,+2,0)\\
60&(61,0,0)&(59,2,0)&(-2,+2,0)\\
62&(63,0,0)&(61,2,0)&(-2,+2,0)\\
64&(65,0,0)&(63,2,0)&(-2,+2,0).
\end{array}
\]
Les propositions $(-2,-1,0)$ et $(2,1,0)$ ne peuvent pas être des différences
d'inertie : leurs sommes valent respectivement $-3$ et $3$, donc elles changeraient
le rang. Leur correction typée est $(-2,+2,0)$ et $(+2,-2,0)$.

\section{Conclusion et frontières}

Au degré $64$, l'extension finie est certifiée par
\[
\In(Q_{\mathrm{pair}})=(33,0,0),\quad
\In(Q_{\mathrm{impair}})=(32,0,0),
\]
\[
\In(L_{\mathrm{pair}})=(32,1,0),\quad
\In(L_{\mathrm{impair}})=(31,1,0).
\]
La tour abstraite des entrelaceurs est $\T$; la covariance par congruence est $\T$;
le transfert ternaire fini est $\T$. Le fil brut testé est $\NT$ entre $62$ et $64$.
Les identités décimales sont $\T$ comme identités de mots, mais leur promotion en
morphisme analytique supplémentaire est ouverte et n'est pas utilisée.

Les gardes finales sont : calcul fini seulement; aucune permutation des entrées brutes;
aucune extrapolation au degré suivant; aucune autorisation CORE7Q\_A2; aucune
revendication RH; complément infini, inertie de Paley--Wiener complète et positivité
globale ouverts.

\end{document}
